% EcoCollect functional test report.
% Build from this folder:  tectonic test-report.tex --outdir ..
\documentclass[11pt]{article}
\usepackage{fontspec}
\usepackage[letterpaper,margin=0.9in]{geometry}
\usepackage{graphicx,longtable,booktabs,array,xcolor,enumitem}
\usepackage{xurl}
\usepackage[colorlinks=true,linkcolor=linkblue,urlcolor=linkblue]{hyperref}

\definecolor{linkblue}{HTML}{1F4E9A}
\definecolor{passgreen}{HTML}{1E6B34}
\definecolor{failred}{HTML}{A32020}
\setlength{\parindent}{0pt}
\setlength{\parskip}{6pt}
\setcounter{tocdepth}{1}
\emergencystretch=3em

\newcommand{\code}[1]{\texttt{#1}}
\newcommand{\passtag}{\textbf{\color{passgreen}Pass}}
\newcommand{\failtag}{\textbf{\color{failred}Fail}}

% ID, scenario, result, setup and steps, expected result, how it was checked
\newcommand{\testcase}[6]{%
  \par\addvspace{10pt}
  \noindent\rule{\linewidth}{0.4pt}\par\nopagebreak
  \phantomsection\label{#1}\textbf{#1\quad #2}\hfill #3\par\nopagebreak
  \begin{description}[style=nextline,leftmargin=1.2em,itemsep=2pt,topsep=2pt,font=\normalfont\itshape]
    \item[Setup and steps] #4
    \item[Expected result] #5
    \item[How it was checked] #6
  \end{description}}

\newcommand{\failreason}[1]{\textbf{\color{failred}Why it failed:} #1\par}

% screenshot name, width as a share of the line, caption
\newcommand{\shot}[3]{%
  \par\addvspace{6pt}
  \noindent\begin{minipage}{\linewidth}\centering
    \fbox{\includegraphics[width=#2\linewidth,height=0.74\textheight,keepaspectratio]{screenshots/#1.png}}\par\smallskip
    \small #3
  \end{minipage}\par\addvspace{6pt}}

% team document, page number, link targets for the cases starting on it
\newcommand{\record}[3]{%
  \par\noindent\begin{minipage}{\linewidth}\centering
    \leavevmode#3\par\nointerlineskip
    \fbox{\includegraphics[page=#2,width=0.94\linewidth,height=0.86\textheight,keepaspectratio]{team-records/#1.pdf}}
  \end{minipage}\par\clearpage}

\begin{document}

\begin{center}
  {\LARGE\bfseries EcoCollect: Functional Test Report}\par
  {\large ENSF 607 --- Software Quality}\par
  November 27, 2026
\end{center}

\tableofcontents

\section{Results}

All testing is manual. The 23 functional requirements have a positive (P) and a negative (N) case each, 46 in total: \textbf{45 pass and 1 fails}. Every case has been run.

\begin{center}
\begin{tabular}{p{0.78\linewidth}r}
\toprule
How the case was run & Passed \\
\midrule
Team manual runs & 34 \\
Runs on the local Compose stack (FR02-P, FR03-N, FR08-P, FR13-N, FR14-N, FR16-P, FR16-N, FR17-N, FR19-P, FR19-N and FR22-P) & 11 \\
\bottomrule
\end{tabular}
\end{center}

\textbf{Failed: \hyperref[FR11-N]{FR11-N}.} After the pickup window ended, the resident page still showed Cancel assignment. The app refused the cancellation and the request did not change. This run used one account as both resident and volunteer.

\subsection*{Where the evidence is}
\begin{itemize}[itemsep=2pt]
  \item \textbf{Team runs:} the team's own documents are included in full in Appendices~\ref{app:A} to~\ref{app:C}. The team ran 38 manual cases and marked all of them as passed. Each document numbers its cases from TC001, so a team case is named by its appendix and number. The case entries use some of the team cases; the rest are extra checks.
  \item \textbf{Screenshots:} Appendix~\ref{app:A} (dated November 26) and Appendix~\ref{app:B} have the team's screenshots. Appendix~\ref{app:C} has none.
  \item \textbf{Local runs:} screenshots appear under FR03-N, FR11-N, FR13-N, FR14-N, FR16-P, FR16-N, FR17-N, FR19-P, FR19-N and FR22-P. FR02-P and FR08-P have none. They are cropped to the part of the page being checked, and personal details are blacked out.
\end{itemize}

\subsection*{Limits}
\begin{itemize}[itemsep=2pt]
  \item The tested commit was not recorded for the team runs or for FR02-P and FR08-P. The other local runs on 2026-11-27 in Chrome, including the failed one, used app code from \code{main} at \code{73700f2}. FR22-P also relies on two team cases.
  \item At that commit the seeded terms were only the version 2 placeholder. \href{https://github.com/dannydufau/ecocollect/pull/511}{PR \#511} later added version 3, which these runs did not test.
  \item A password visible in the TC001 screenshot of Appendix~\ref{app:A} is blacked out.
  \item This report covers the functional requirements. Acceptance and system tests are kept separately in \code{cases/acceptance.md}.
\end{itemize}

\section{Coverage by requirement}
\begin{longtable}{lp{0.62\linewidth}cc}
\toprule
ID & Functional requirement & P & N \\
\midrule
\endhead
\bottomrule
\endfoot
\hyperref[FR01]{FR01} & Account Registration & Pass & Pass \\
\hyperref[FR02]{FR02} & User Login, Logout, and Role Selection & Pass & Pass \\
\hyperref[FR03]{FR03} & Additional Role Activation & Pass & Pass \\
\hyperref[FR04]{FR04} & Profile Management & Pass & Pass \\
\hyperref[FR05]{FR05} & Public User Profiles & Pass & Pass \\
\hyperref[FR06]{FR06} & Create Pickup Requests & Pass & Pass \\
\hyperref[FR07]{FR07} & Pickup Scheduling & Pass & Pass \\
\hyperref[FR08]{FR08} & Review and Confirm Pickup Requests & Pass & Pass \\
\hyperref[FR09]{FR09} & View Resident Pickup Requests & Pass & Pass \\
\hyperref[FR10]{FR10} & Edit Open Pickup Requests & Pass & Pass \\
\hyperref[FR11]{FR11} & Cancel Pickup Requests & Pass & \textbf{\color{failred}Fail} \\
\hyperref[FR12]{FR12} & Volunteer Request Discovery & Pass & Pass \\
\hyperref[FR13]{FR13} & Accept Pickup Requests & Pass & Pass \\
\hyperref[FR14]{FR14} & Cancel Volunteer Assignment & Pass & Pass \\
\hyperref[FR15]{FR15} & Confirm Item Collection & Pass & Pass \\
\hyperref[FR16]{FR16} & Report Missed Pickup & Pass & Pass \\
\hyperref[FR17]{FR17} & Record Facility Delivery & Pass & Pass \\
\hyperref[FR18]{FR18} & EcoPoints Calculation and Awarding & Pass & Pass \\
\hyperref[FR19]{FR19} & In-App Notifications and Push Notifications & Pass & Pass \\
\hyperref[FR20]{FR20} & Automatic Request Expiration & Pass & Pass \\
\hyperref[FR21]{FR21} & Volunteer Pickup History & Pass & Pass \\
\hyperref[FR22]{FR22} & Terms and Conditions, including Safety Guidelines & Pass & Pass \\
\hyperref[FR23]{FR23} & Public Community Impact Statistics & Pass & Pass \\
\end{longtable}

\section{How the cases are run}

In each case ID, P is a positive case that checks the expected result of valid use, and N is a negative case that checks a rejection or a constraint. All cases are run by hand.

\begin{enumerate}[itemsep=2pt]
  \item Start the app with Docker Compose as in the project README, using Firebase test accounts.
  \item Follow \emph{Setup and steps} and compare what happens with \emph{Expected result}.
  \item Record Pass, or Fail with the reason. Note the date and commit, and save screenshots.
\end{enumerate}

For FR11-N and FR16, accept a request before its window starts. Check FR16-N before the window ends, then FR11-N and FR16-P after it ends, in that order.

For FR19, set the Web Push keys in \code{.env} (see \code{.env.example}), use a certificate the browser trusts, and use normal Chrome profiles. Push does not work in Incognito.

\subsection*{Database and event flows}
These cases check behaviour that depends on events between the services and on stored data:
\begin{itemize}[itemsep=2pt]
  \item \textbf{FR20-P and FR20-N:} a request expires, or stays Open, with no user action.
  \item \textbf{FR18-P and FR23-P:} after a delivery, the EcoPoints awarded and the community statistics change.
  \item \textbf{FR19-P:} a request event creates an in-app notification and a push.
  \item \textbf{FR19-N:} checks, with the notification tables, that a deactivated volunteer gets nothing and that cancelling a request expires the old notification.
\end{itemize}

\clearpage
\section{Test cases}
\filbreak
\subsection{FR01: Account Registration}\label{FR01}
\testcase{FR01-P}{Register as a resident}{\passtag}
  {While signed out, open Register and choose Resident. Fill in the name, public name, username, a new email, date of birth, password, street address and postal code. Tick the terms box and choose Create account.}
  {The account is created and signed in. The Continue as page offers the Resident role.}
  {Manual --- team run, \hyperref[app:A]{FR01--05 and FR11--15 cases (Appendix A)}, \hyperlink{tc:A01}{TC001}.}

\filbreak
\testcase{FR01-N}{Reject a taken username}{\passtag}
  {With an account already using a username, register again with that username and a different new email. Tick the terms box and choose Create account.}
  {An error next to Username says the username is already taken. The user stays on Register and no second account is created.}
  {Manual --- team run, \hyperref[app:A]{FR01--05 and FR11--15 cases (Appendix A)}, \hyperlink{tc:A02}{TC002}.}

\filbreak
\subsection{FR02: User Login, Logout, and Role Selection}\label{FR02}
\testcase{FR02-P}{Sign in, choose a role, and sign out}{\passtag}
  {On a scratch stack with the frontend and Firebase Auth configured, sign in with a dual-role test account. Check that Resident and Volunteer choices appear, enter one role, and log out. Try opening the protected page again.}
  {The app offers both role choices after sign-in. The selected role opens its view, and logout ends the session so the protected page requires sign-in again.}
  {Manual --- ran on the local Compose stack with Firebase Auth: signed in with a dual-role account, checked both role choices on \code{/choose-role}, opened \code{/resident}, logged out, and checked that \code{/resident} redirected to \code{/login}.}

\filbreak
\testcase{FR02-N}{Reject a wrong password}{\passtag}
  {While signed out, enter a registered email with a wrong password and choose Sign in.}
  {The page says “Error. Incorrect email or password.” The user stays signed out on the sign-in page.}
  {Manual --- team run, \hyperref[app:A]{FR01--05 and FR11--15 cases (Appendix A)}, \hyperlink{tc:A04}{TC004}.}

\filbreak
\subsection{FR03: Additional Role Activation}\label{FR03}
\testcase{FR03-P}{Activate the second role}{\passtag}
  {With a resident-only account, open the profile page and choose Become a volunteer. Add a preferred postal code, tick the terms box and choose Activate role.}
  {The profile shows both roles for the same account, and the Continue as page offers Resident and Volunteer.}
  {Manual --- team run, \hyperref[app:A]{FR01--05 and FR11--15 cases (Appendix A)}, \hyperlink{tc:A05}{TC005}.}

\filbreak
\testcase{FR03-N}{Do not offer a role the account already has}{\passtag}
  {With an account that has both roles, open the profile page.}
  {There is no Add a role section, so neither role can be added again.}
  {Manual --- run in Chrome on 2026-11-27 on the local Compose stack (app code from \code{main} at \code{73700f2}), using one account with both roles. After the Volunteer role was added and the profile reloaded, there was no Add a role section.}
\shot{fr03-n}{0.5}{FR03-N: Right after the Volunteer role was added: both roles are listed and there is no Add a role section.}

\filbreak
\subsection{FR04: Profile Management}\label{FR04}
\testcase{FR04-P}{View the profile and change a setting}{\passtag}
  {Signed in, open the profile page, edit the public name and save. Then open the account's public profile.}
  {The profile shows the new public name and confirms it was saved. The public profile shows the new name.}
  {Manual --- team run, \hyperref[app:A]{FR01--05 and FR11--15 cases (Appendix A)}, \hyperlink{tc:A07}{TC007}.}

\filbreak
\testcase{FR04-N}{Block deactivation while a request is active}{\passtag}
  {As a resident who owns an Open request, choose Deactivate account on the profile page and confirm.}
  {The page says the account can't be deactivated while requests or pickups are in progress. The account stays active.}
  {Manual --- team run, \hyperref[app:A]{FR01--05 and FR11--15 cases (Appendix A)}, \hyperlink{tc:A08}{TC008}.}

\filbreak
\subsection{FR05: Public User Profiles}\label{FR05}
\testcase{FR05-P}{Show only public participation data}{\passtag}
  {While signed out, open the public profile link of an account that has both roles.}
  {The page shows the public name, roles, member-since date and EcoPoints. It does not show the email, birth date, address or any pickup request details.}
  {Manual --- team run, \hyperref[app:A]{FR01--05 and FR11--15 cases (Appendix A)}, \hyperlink{tc:A09}{TC009}.}

\filbreak
\testcase{FR05-N}{Show nothing for an unknown profile}{\passtag}
  {While signed out, open a public profile link with a made-up ID.}
  {The page shows “Profile not found” and no account data.}
  {Manual --- team run, \hyperref[app:A]{FR01--05 and FR11--15 cases (Appendix A)}, \hyperlink{tc:A10}{TC010}.}

\filbreak
\subsection{FR06: Create Pickup Requests}\label{FR06}
\testcase{FR06-P}{Create a pickup request}{\passtag}
  {As a resident, choose Computers/Laptops with quantity 2 and one window for tomorrow. Review, tick the terms box and choose Confirm and send.}
  {The form estimates 5.00 kg. My requests lists the new request as Open.}
  {Manual --- team run, \hyperref[app:B]{FR06--10 cases (Appendix B)}, \hyperlink{tc:B01}{TC001}.}

\filbreak
\testcase{FR06-N}{Reject an unsupported photo}{\passtag}
  {On a new request, choose Televisions with quantity 1 and pick a GIF file as the photo.}
  {The photo is not uploaded and the form asks for a JPEG, PNG or WebP photo.}
  {Manual --- team run, \hyperref[app:B]{FR06--10 cases (Appendix B)}, \hyperlink{tc:B02}{TC002}.}

\filbreak
\subsection{FR07: Pickup Scheduling}\label{FR07}
\testcase{FR07-P}{Offer three pickup windows}{\passtag}
  {As a resident, start a request for one phone, add three one-hour windows for tomorrow, review, accept the terms and send. Open the request.}
  {The request is Open and shows all three windows.}
  {Manual --- team run, \hyperref[app:B]{FR06--10 cases (Appendix B)}, \hyperlink{tc:B03}{TC003}.}

\filbreak
\testcase{FR07-N}{Require at least one pickup window}{\passtag}
  {On a new request, choose one phone, remove the pickup window and choose Review.}
  {The review page does not open and the form says to choose one to three pickup windows. No request is created.}
  {Manual --- team run, \hyperref[app:B]{FR06--10 cases (Appendix B)}, \hyperlink{tc:B04}{TC004}.}

\filbreak
\subsection{FR08: Review and Confirm Pickup Requests}\label{FR08}
\testcase{FR08-P}{Correct a pickup request before confirming}{\passtag}
  {On the integrated scratch frontend, create a resident draft, open the review page, return to the form, correct the quantity, and review again. Acknowledge the current terms and confirm.}
  {The review displays the correction. Only the final confirmation creates the request, using the corrected quantity.}
  {Manual --- ran on the local Compose stack with Firebase Auth: drafted \code{1 × Computers/Laptops}, opened Review (\code{2.50 kg}), went back to change quantity to \code{3}, checked \code{7.50 kg} on Review, accepted terms, confirmed, and verified the created \code{Open} request.}

\filbreak
\testcase{FR08-N}{Create nothing before confirming}{\passtag}
  {Fill in a new request and choose Review. Open My requests without sending it.}
  {My requests does not list a new request.}
  {Manual --- team run, \hyperref[app:B]{FR06--10 cases (Appendix B)}, \hyperlink{tc:B06}{TC006}.}

\filbreak
\subsection{FR09: View Resident Pickup Requests}\label{FR09}
\testcase{FR09-P}{Filter my requests by status}{\passtag}
  {As a resident with an Open, a Cancelled and an Assigned request, choose the Assigned filter on My requests.}
  {Only the Assigned request is listed.}
  {Manual --- team run, \hyperref[app:B]{FR06--10 cases (Appendix B)}, \hyperlink{tc:B07}{TC007}.}

\filbreak
\testcase{FR09-N}{Filter with no matching request}{\passtag}
  {With the same three requests, choose the Delivered filter.}
  {No request is listed. The page says there are no delivered requests and offers Show all requests.}
  {Manual --- team run, \hyperref[app:B]{FR06--10 cases (Appendix B)}, \hyperlink{tc:B08}{TC008}.}

\filbreak
\subsection{FR10: Edit Open Pickup Requests}\label{FR10}
\testcase{FR10-P}{Edit the quantity of an Open request}{\passtag}
  {Edit an Open request for two laptops, change the quantity to 5 and save.}
  {The estimated weight becomes 12.50 kg and the request stays Open.}
  {Manual --- team run, \hyperref[app:B]{FR06--10 cases (Appendix B)}, \hyperlink{tc:B09}{TC009}.}

\filbreak
\testcase{FR10-N}{Reject a quantity of zero}{\passtag}
  {Edit an Open request, change the quantity to 0 and try to save.}
  {The form asks for a whole number of 1 or more. The request is not changed.}
  {Manual --- team run, \hyperref[app:B]{FR06--10 cases (Appendix B)}, \hyperlink{tc:B10}{TC010}.}

\filbreak
\subsection{FR11: Cancel Pickup Requests}\label{FR11}
\testcase{FR11-P}{Cancel an Open pickup request}{\passtag}
  {On an Open request, choose Cancel request, give a reason and confirm.}
  {The request status changes to Cancelled.}
  {Manual --- team run, \hyperref[app:A]{FR01--05 and FR11--15 cases (Appendix A)}, \hyperlink{tc:A11}{TC011}.}

\filbreak
\testcase{FR11-N}{No cancelling after the selected window ends}{\failtag}
  {Take a request that was accepted and not collected. After its selected window ends, open it as the resident.}
  {There is no option to cancel the request or the assignment. Confirm collected and Report missed pickup are offered.}
  {Manual --- run in Chrome on 2026-11-27 on the local Compose stack (app code from \code{main} at \code{73700f2}), using one account with both roles. After the window ended, Cancel request was gone and Confirm collected and Report missed pickup were offered, but Cancel assignment was still shown. Confirming it showed “The pickup window has ended” and the request stayed Assigned. From the code, the button appears because the same account may still cancel as the volunteer, while the resident page cancels as the resident. Not tried with separate accounts.}
\failreason{Cancel assignment was still offered after the window ended. The cancellation itself was refused.}
\shot{fr11-n-1}{0.85}{FR11-N: After the window ended, Cancel assignment is still shown next to Report missed pickup.}
\shot{fr11-n-2}{0.52}{FR11-N: Confirming it is refused: “The pickup window has ended.”}

\filbreak
\subsection{FR12: Volunteer Request Discovery}\label{FR12}
\testcase{FR12-P}{List requests in the saved postal code}{\passtag}
  {As a volunteer with one saved postal code and an Open request in that area, open Available requests.}
  {The filter shows the saved code and the request in that area is listed.}
  {Manual --- team run, \hyperref[app:A]{FR01--05 and FR11--15 cases (Appendix A)}, \hyperlink{tc:A13}{TC013}.}

\filbreak
\testcase{FR12-N}{Filter on a code with no requests}{\passtag}
  {Replace the filter with a postal code that has no Open requests, then choose Show all areas.}
  {The filtered list says there are no open requests in these areas. Show all areas lists Open requests again.}
  {Manual --- team run, \hyperref[app:A]{FR01--05 and FR11--15 cases (Appendix A)}, \hyperlink{tc:A14}{TC014}.}

\filbreak
\subsection{FR13: Accept Pickup Requests}\label{FR13}
\testcase{FR13-P}{Accept a request}{\passtag}
  {Open a request with future windows, choose Accept request, pick a window, tick the terms box and confirm.}
  {The request is Assigned with the chosen window, and the volunteer can see the pickup address.}
  {Manual --- team run, \hyperref[app:A]{FR01--05 and FR11--15 cases (Appendix A)}, \hyperlink{tc:A15}{TC015}. Extra check in the same document, \hyperlink{tc:A16}{TC016}: an account with both roles accepted its own request.}

\filbreak
\testcase{FR13-N}{Refuse a second volunteer}{\passtag}
  {Two volunteers open the same Open request in separate browser profiles. Volunteer A accepts it, then volunteer B tries to accept it.}
  {B cannot accept it: the page stops offering it, or says “Another volunteer accepted this request first.” The request stays assigned to A.}
  {Manual --- run in Chrome on 2026-11-27 on the local Compose stack (app code from \code{main} at \code{73700f2}), with two volunteer accounts in a normal and an Incognito window. The tester confirmed the second window had the Open request on screen before the first volunteer accepted it. That window then showed “This request is no longer available”. The database showed one assignment, to the first volunteer.}
\shot{fr13-n}{0.95}{FR13-N: The second volunteer’s window after the first volunteer accepted.}

\filbreak
\subsection{FR14: Cancel Volunteer Assignment}\label{FR14}
\testcase{FR14-P}{Cancel an assignment before collection}{\passtag}
  {On an Assigned request with future windows, the volunteer chooses Cancel assignment, gives a reason and confirms.}
  {The request goes back to Open and is listed in Available requests again.}
  {Manual --- team run, \hyperref[app:A]{FR01--05 and FR11--15 cases (Appendix A)}, \hyperlink{tc:A17}{TC017}.}

\filbreak
\testcase{FR14-N}{No cancelling after collection}{\passtag}
  {After the resident confirms collection, open the request as the assigned volunteer.}
  {There is no option to cancel the assignment. The request stays Collected.}
  {Manual --- run in Chrome on 2026-11-27 on the local Compose stack (app code from \code{main} at \code{73700f2}), using one account with both roles. After collection was confirmed, the volunteer page offered Record delivery and no way to cancel. The request stayed Collected.}
\shot{fr14-n}{0.72}{FR14-N: The Collected request offers Record delivery and no way to cancel.}

\filbreak
\subsection{FR15: Confirm Item Collection}\label{FR15}
\testcase{FR15-P}{The owner confirms collection}{\passtag}
  {As the resident who owns an Assigned request, choose Confirm collected.}
  {The request is Collected and its history shows the change.}
  {Manual --- team run, \hyperref[app:A]{FR01--05 and FR11--15 cases (Appendix A)}, \hyperlink{tc:A19}{TC019}.}

\filbreak
\testcase{FR15-N}{The volunteer cannot confirm collection}{\passtag}
  {As the assigned volunteer, open the request.}
  {The page has no Confirm collected button and the request stays Assigned.}
  {Manual --- team run, \hyperref[app:A]{FR01--05 and FR11--15 cases (Appendix A)}, \hyperlink{tc:A20}{TC020}.}

\filbreak
\subsection{FR16: Report Missed Pickup}\label{FR16}
\testcase{FR16-P}{Report a missed pickup and reschedule}{\passtag}
  {After the selected window of an accepted, uncollected request ends, the resident chooses Report missed pickup, gives a reason, adds a new window, accepts the terms and confirms.}
  {The old request is Missed and a new linked Open request with the new window is shown.}
  {Manual --- run in Chrome on 2026-11-27 on the local Compose stack (app code from \code{main} at \code{73700f2}), using one account with both roles. After the window ended, reported the pickup missed with a reason, added a window for the next day and accepted the terms. The old request became Missed with the reason in its history, and a linked Open request showed the new window.}
\shot{fr16-p-1}{0.52}{FR16-P: The missed pickup form with a reason and a new window, before the terms box was ticked. The Report missed pickup button before this is visible in the first FR11-N screenshot.}
\shot{fr16-p-2}{0.85}{FR16-P: The new Open request, linked to the earlier one. There is no screenshot of the earlier request showing Missed.}

\filbreak
\testcase{FR16-N}{No early missed-pickup report}{\passtag}
  {Before the selected window of an accepted request ends, open it as the resident.}
  {Report missed pickup is not offered. The request stays Assigned.}
  {Manual --- run in Chrome on 2026-11-27 on the local Compose stack (app code from \code{main} at \code{73700f2}), using one account with both roles. Before the accepted window started, the resident page offered no Report missed pickup. The request stayed Assigned.}
\shot{fr16-n}{0.85}{FR16-N: Before the window starts, Report missed pickup is not offered.}

\filbreak
\subsection{FR17: Record Facility Delivery}\label{FR17}
\testcase{FR17-P}{Record a delivery}{\passtag}
  {As the volunteer with a Collected request, record the delivery: choose a facility and add a delivery photo.}
  {The request is Delivered.}
  {Manual --- team run, \hyperref[app:C]{FR20--23 cases (Appendix C)}, \hyperlink{tc:C07}{TC007}.}

\filbreak
\testcase{FR17-N}{No delivery before collection}{\passtag}
  {Before the resident confirms collection, open the request as the assigned volunteer.}
  {There is no option to record a delivery. The request stays Assigned.}
  {Manual --- run in Chrome on 2026-11-27 on the local Compose stack (app code from \code{main} at \code{73700f2}), using one account with both roles. While the request was Assigned, the volunteer page offered Cancel assignment and no Record delivery.}
\shot{fr17-n}{0.85}{FR17-N: While Assigned, the volunteer page offers Cancel assignment and no Record delivery.}

\filbreak
\subsection{FR18: EcoPoints Calculation and Awarding}\label{FR18}
\testcase{FR18-P}{Award points for a delivery}{\passtag}
  {Note EcoPoints awarded on the Community impact page, record the delivery of one laptop (2.50 kg), then refresh.}
  {EcoPoints awarded goes up by 19.50, the total for this delivery (13 for the volunteer and 6.50 for the resident; the page shows only the total).}
  {Manual --- team run, \hyperref[app:C]{FR20--23 cases (Appendix C)}, \hyperlink{tc:C07}{TC007}.}

\filbreak
\testcase{FR18-N}{No points without a delivery}{\passtag}
  {Note EcoPoints awarded. Create a request and cancel it, confirm collection of another request, then refresh.}
  {EcoPoints awarded does not change.}
  {Manual --- team run, \hyperref[app:C]{FR20--23 cases (Appendix C)}, \hyperlink{tc:C08}{TC008}.}

\filbreak
\subsection{FR19: In-App Notifications and Push Notifications}\label{FR19}
\testcase{FR19-P}{Notify a matched volunteer and mark the notification read}{\passtag}
  {Sign in as a volunteer with a preferred postal code and allow notifications. In a separate browser profile, submit a request as a resident at an address with that postal code. Check the volunteer’s notification list and device alert, mark the in-app notification as read and reload.}
  {The volunteer gets the new-request notification in the app and by push. It is still read after the reload.}
  {Manual --- run in Chrome on 2026-11-27 on the local Compose stack (app code from \code{main} at \code{73700f2}) with the local certificate trusted. With push on for the volunteer, the resident submitted a request for one monitor from a separate browser session. The volunteer’s list showed “New pickup near you”, the database recorded the push at 7:51 p.m. and the tester reported the Mac alert. The tester then opened the notification and reloaded; it stayed read, and the database showed it as read.}
\shot{fr19-p-1}{0.9}{FR19-P: Push notifications on in the volunteer’s browser, before the request was sent.}
\shot{fr19-p-2}{0.9}{FR19-P: The volunteer’s list with the new matching request.}
\shot{fr19-after}{0.9}{FR19-P: Later in the same session, the monitor notification still shows as read. The Mac alert was not captured.}

\filbreak
\testcase{FR19-N}{Skip a deactivated volunteer and clear an outdated notification}{\passtag}
  {Use volunteers A and B with the same preferred postal code. B allows notifications, then deactivates the account. A resident submits a request in that postal code. Check A, and check the Notification Service tables for both accounts. Without opening or reading A’s notification, cancel the request, then check A’s list and device.}
  {Only A gets the notification. The database has no new stored notification or push subscription for B. After the cancellation it leaves A’s list. On the device, the unread alert is closed or replaced by a silent “This pickup is no longer available”.}
  {Manual --- run in Chrome on 2026-11-27 on the local Compose stack (app code from \code{main} at \code{73700f2}). Volunteer B turned on push and then deactivated the account. After the resident submitted a request for one television, volunteer A’s list showed “New pickup near you” and the database recorded the push to A. B had no stored notification or push subscription. Without A opening it, the resident cancelled the request. The database marked A’s unread notification expired, and it was gone from A’s list. The tester reported seeing the Mac alert before the cancellation and that it was gone afterwards.}
\shot{fr19-n-1}{0.62}{FR19-N: Volunteer B with push on, before deactivating.}
\shot{fr19-n-2}{0.9}{FR19-N: Volunteer A’s list after the television request was sent.}
\shot{fr19-after}{0.9}{FR19-N: Volunteer A’s list after the resident cancelled: the television notification is gone. The Mac alert was not captured.}

\filbreak
\subsection{FR20: Automatic Request Expiration}\label{FR20}
\testcase{FR20-P}{Expire when the last window starts}{\passtag}
  {Create a request with one window at the next full hour. Wait until it starts and refresh. Then check Available requests as a volunteer.}
  {The request changes from Open to Expired without any user action, and volunteers no longer see it.}
  {Manual --- team run, \hyperref[app:C]{FR20--23 cases (Appendix C)}, \hyperlink{tc:C01}{TC001}.}

\filbreak
\testcase{FR20-N}{Stay Open while a later window remains}{\passtag}
  {Create a request with two windows an hour apart. Wait until the first one starts, then open it as a volunteer.}
  {The request is still Open, and only the later window is offered.}
  {Manual --- team run, \hyperref[app:C]{FR20--23 cases (Appendix C)}, \hyperlink{tc:C02}{TC002}.}

\filbreak
\subsection{FR21: Volunteer Pickup History}\label{FR21}
\testcase{FR21-P}{View pickup history}{\passtag}
  {As a volunteer with Active, Collected and Delivered assignments, open My assignments, then choose the Delivered filter.}
  {All three show their item, quantity, status, window and acceptance time. The filter shows only the Delivered one.}
  {Manual --- team run, \hyperref[app:C]{FR20--23 cases (Appendix C)}, \hyperlink{tc:C03}{TC003}.}

\filbreak
\testcase{FR21-N}{Hide the details of a cancelled request}{\passtag}
  {After the resident cancels an accepted request, the volunteer opens My assignments, then the request's link.}
  {The row says Cancelled by resident without the category, quantity or postal code. The link says the request is no longer available.}
  {Manual --- team run, \hyperref[app:C]{FR20--23 cases (Appendix C)}, \hyperlink{tc:C04}{TC004}.}

\filbreak
\subsection{FR22: Terms and Conditions, including Safety Guidelines}\label{FR22}
\testcase{FR22-P}{Require current terms in all four flows}{\passtag}
  {Open registration, role activation, request review and pickup acceptance. In each, check that the current terms are shown and try to submit with the terms box unticked. Then tick the box and accept a pickup.}
  {Each flow shows the current terms with an acknowledgement box, and nothing is submitted until it is ticked. After it is ticked, the pickup is accepted.}
  {Manual --- run in Chrome on 2026-11-27 on the local Compose stack (app code from \code{main} at \code{73700f2}), using one account with both roles. Activate role, Confirm and send, and Accept request stayed disabled until the box was ticked, and the pickup was then accepted. Registration: team run, \hyperref[app:C]{FR20--23 cases (Appendix C)}, \hyperlink{tc:C06}{TC006}, with the terms shown in \hyperref[app:A]{FR01--05 and FR11--15 cases (Appendix A)}, \hyperlink{tc:A01}{TC001}. At that commit the seeded terms were only the version 2 placeholder, so these three screens showed no safety guidelines text. \href{https://github.com/dannydufau/ecocollect/pull/511}{PR \#511} later added version 3, which this run did not test.}
\shot{fr22-p-1}{0.55}{FR22-P: Role activation: Activate role is disabled.}
\shot{fr22-p-2}{0.46}{FR22-P: Request review: Confirm and send is disabled.}
\shot{fr22-p-3}{0.46}{FR22-P: Pickup acceptance: Accept request is disabled.}

\filbreak
\testcase{FR22-N}{Block registration without accepting the terms}{\passtag}
  {On Register, fill every field with valid resident details, leave the terms box unticked and try to choose Create account.}
  {Create account stays disabled. The user stays on Register and no account is created.}
  {Manual --- team run, \hyperref[app:C]{FR20--23 cases (Appendix C)}, \hyperlink{tc:C06}{TC006}.}

\filbreak
\subsection{FR23: Public Community Impact Statistics}\label{FR23}
\testcase{FR23-P}{Count a new delivery}{\passtag}
  {Note the Community impact values, record the delivery of one laptop (2.50 kg), then refresh the page.}
  {Deliveries and items go up by 1, weight by 2.50 kg and EcoPoints awarded by 19.50. The Computers row includes it.}
  {Manual --- team run, \hyperref[app:C]{FR20--23 cases (Appendix C)}, \hyperlink{tc:C07}{TC007}.}

\filbreak
\testcase{FR23-N}{Leave out requests that were not delivered}{\passtag}
  {Note the Community impact values. Cancel a new request, confirm collection of another, then refresh the page.}
  {The values do not change.}
  {Manual --- team run, \hyperref[app:C]{FR20--23 cases (Appendix C)}, \hyperlink{tc:C08}{TC008}.}



\appendix
\clearpage
\section{Team record: FR01--05 and FR11--15}\label{app:A}
\record{fr01-05-fr11-15}{1}{\hypertarget{tc:A01}{}}
\record{fr01-05-fr11-15}{2}{}
\record{fr01-05-fr11-15}{3}{\hypertarget{tc:A02}{}}
\record{fr01-05-fr11-15}{4}{\hypertarget{tc:A03}{}}
\record{fr01-05-fr11-15}{5}{}
\record{fr01-05-fr11-15}{6}{\hypertarget{tc:A04}{}}
\record{fr01-05-fr11-15}{7}{}
\record{fr01-05-fr11-15}{8}{\hypertarget{tc:A05}{}}
\record{fr01-05-fr11-15}{9}{\hypertarget{tc:A06}{}}
\record{fr01-05-fr11-15}{10}{\hypertarget{tc:A07}{}}
\record{fr01-05-fr11-15}{11}{\hypertarget{tc:A08}{}}
\record{fr01-05-fr11-15}{12}{\hypertarget{tc:A09}{}}
\record{fr01-05-fr11-15}{13}{\hypertarget{tc:A10}{}}
\record{fr01-05-fr11-15}{14}{\hypertarget{tc:A11}{}}
\record{fr01-05-fr11-15}{15}{}
\record{fr01-05-fr11-15}{16}{\hypertarget{tc:A12}{}}
\record{fr01-05-fr11-15}{17}{\hypertarget{tc:A13}{}}
\record{fr01-05-fr11-15}{18}{\hypertarget{tc:A14}{}}
\record{fr01-05-fr11-15}{19}{\hypertarget{tc:A15}{}}
\record{fr01-05-fr11-15}{20}{}
\record{fr01-05-fr11-15}{21}{\hypertarget{tc:A16}{}}
\record{fr01-05-fr11-15}{22}{\hypertarget{tc:A17}{}}
\record{fr01-05-fr11-15}{23}{}
\record{fr01-05-fr11-15}{24}{\hypertarget{tc:A18}{}}
\record{fr01-05-fr11-15}{25}{\hypertarget{tc:A19}{}}
\record{fr01-05-fr11-15}{26}{\hypertarget{tc:A20}{}}
\record{fr01-05-fr11-15}{27}{}
\clearpage
\section{Team record: FR06--10}\label{app:B}
\record{fr06-10}{1}{\hypertarget{tc:B01}{}}
\record{fr06-10}{2}{}
\record{fr06-10}{3}{\hypertarget{tc:B02}{}}
\record{fr06-10}{4}{\hypertarget{tc:B03}{}}
\record{fr06-10}{5}{}
\record{fr06-10}{6}{\hypertarget{tc:B04}{}}
\record{fr06-10}{7}{\hypertarget{tc:B05}{}}
\record{fr06-10}{8}{\hypertarget{tc:B06}{}}
\record{fr06-10}{9}{\hypertarget{tc:B07}{}}
\record{fr06-10}{10}{\hypertarget{tc:B08}{}}
\record{fr06-10}{11}{\hypertarget{tc:B09}{}}
\record{fr06-10}{12}{}
\record{fr06-10}{13}{\hypertarget{tc:B10}{}}
\record{fr06-10}{14}{}
\clearpage
\section{Team record: FR20--23}\label{app:C}
\record{fr20-23}{1}{\hypertarget{tc:C01}{}\hypertarget{tc:C02}{}}
\record{fr20-23}{2}{\hypertarget{tc:C03}{}\hypertarget{tc:C04}{}}
\record{fr20-23}{3}{\hypertarget{tc:C05}{}\hypertarget{tc:C06}{}}
\record{fr20-23}{4}{\hypertarget{tc:C07}{}\hypertarget{tc:C08}{}}
\record{fr20-23}{5}{}

\end{document}
